% =====================================================================
%  Agent Smith:
%  A Specification Language for Instantiating Purpose-Directed Agents,
%  and a Compiler That Instantiates Them Honouring Four Invariants
%
%  Self-contained. States its own objects; every construct of the
%  language is given a meaning as a finite weighted graph, a bounded
%  convex program, or a construction on one, and the compiler's
%  guarantee is proved from those meanings. The agent's domain mechanics
%  are NOT part of the language: they are supplied by the specifier
%  through the declared drive and scenes. Agent Smith specifies and
%  instantiates; it does not smith.
% =====================================================================
\documentclass[11pt,twocolumn,a4paper]{article}

\usepackage[utf8]{inputenc}
\usepackage[T1]{fontenc}
\usepackage{lmodern}
\usepackage[margin=1in]{geometry}
\usepackage{amsmath,amssymb,amsthm,mathtools}
\usepackage{stmaryrd}
\usepackage{thmtools}
\usepackage{enumitem}
\usepackage{microtype}
\usepackage{booktabs}
\usepackage{graphicx}
\usepackage{listings}
\usepackage{hyperref}
\usepackage{cleveref}
\usepackage{xcolor}

\hypersetup{
  colorlinks=true, linkcolor=blue!50!black,
  citecolor=green!50!black, urlcolor=blue!60!black
}

% ---- theorem environments ----
\declaretheoremstyle[
  spaceabove=6pt, spacebelow=6pt,
  headfont=\bfseries, notefont=\mdseries, notebraces={(}{)},
  bodyfont=\itshape, postheadspace=1em
]{thmstyle}
\declaretheoremstyle[
  spaceabove=6pt, spacebelow=6pt,
  headfont=\bfseries, notefont=\mdseries, notebraces={(}{)},
  bodyfont=\normalfont, postheadspace=1em
]{defstyle}

\declaretheorem[style=thmstyle, name=Theorem, numberwithin=section]{theorem}
\declaretheorem[style=thmstyle, name=Proposition, sibling=theorem]{proposition}
\declaretheorem[style=thmstyle, name=Lemma, sibling=theorem]{lemma}
\declaretheorem[style=thmstyle, name=Corollary, sibling=theorem]{corollary}
\declaretheorem[style=defstyle, name=Definition, sibling=theorem]{definition}
\declaretheorem[style=defstyle, name=Rule, sibling=theorem]{trule}
\declaretheorem[style=defstyle, name=Remark, sibling=theorem]{remark}
\declaretheorem[style=defstyle, name=Example, sibling=theorem]{example}
\declaretheorem[style=defstyle, name=Principle, sibling=theorem]{principle}

\crefname{trule}{Rule}{Rules}
\Crefname{trule}{Rule}{Rules}

% ---- notation ----
\newcommand{\R}{\mathbb{R}}
\newcommand{\N}{\mathbb{N}}
\newcommand{\Agent}{\mathsf{A}}
\newcommand{\Self}{\Gamma}
\newcommand{\V}{V}
\newcommand{\E}{E}
\newcommand{\wt}{w}
\newcommand{\floorc}{\beta}
\newcommand{\cut}{\partial}
\newcommand{\Char}{\chi}
\newcommand{\drive}{\Phi}
\newcommand{\Purp}{x^{\star}}
\newcommand{\att}{\alpha}
\newcommand{\Scene}{\mathsf{S}}
\newcommand{\scenes}{\mathcal{S}}
\newcommand{\Soc}{\Sigma}
\newcommand{\Out}{\mathcal{O}}
\newcommand{\spec}{\mathsf{s}}         % a specification (a script)
\newcommand{\Specs}{\mathbb{S}}        % the language of specifications
\newcommand{\compile}{\mathcal{C}}     % the compiler
\newcommand{\coh}{\kappa}              % coherence
\newcommand{\dist}{d}
\newcommand{\abs}[1]{\lvert#1\rvert}
\newcommand{\norm}[1]{\lVert#1\rVert}
\newcommand{\judge}{\vdash}
\newcommand{\lowers}{\rightsquigarrow} % lowering relation

% ---- listing style for Agent Smith scripts ----
\definecolor{kw}{RGB}{30,64,120}
\definecolor{cm}{RGB}{90,110,90}
\definecolor{st}{RGB}{140,60,40}
\lstdefinelanguage{AgentSmith}{
  morekeywords={agent,drive,purpose,scene,scenes,society,self,parts,
    separations,budget,floor,coherence,phase,outcome,gain,minimise,
    maximise,spawn,snapshot,retire,configure,tie,couple,attends,
    serves,as,when,over,in,with,keeps,reach},
  sensitive=true,
  morecomment=[l]{//},
  morestring=[b]",
}
\lstset{
  language=AgentSmith,
  basicstyle=\ttfamily\footnotesize,
  keywordstyle=\color{kw}\bfseries,
  commentstyle=\color{cm}\itshape,
  stringstyle=\color{st},
  columns=fullflexible,
  keepspaces=true,
  showstringspaces=false,
  breaklines=true,
  frame=single,
  framesep=4pt,
  rulecolor=\color{black!25},
  xleftmargin=2pt,
}

\title{\textbf{Agent Smith:\\
A Specification Language for Instantiating Purpose-Directed Agents,\\
with a Compiler That Instantiates Them Honouring Four Invariants}}

\author{Kundai Farai Sachikonye\\[2pt]
\small AIMe Registry for Artificial Intelligence\\
\small \texttt{kundai.sachikonye@bitspark.com}}
\date{\today}

\begin{document}
\maketitle

\begin{abstract}
\noindent
A game designer specifies a non-player character not by writing its every
motion but by saying what it \emph{is} and what it is \emph{for}: an iron
smith who forges, contracts, and sources ore; a baker who bakes because he
must. We give a language that lets any user, application, or module specify
an agent in exactly this way---by declaring what the agent is for, not how
each act is performed---and a compiler that turns the specification into a
running agent. The language is \textbf{Agent Smith}; every script written in
it \emph{is} an Agent Smith, a complete specification of one agent. We fix
the language's objects as those of the underlying theory of purpose-directed
agents: a self-graph carrying a conserved identity, a purpose that is the
attractor of a declared drive, a finite family of scenes the agent may act
in, a coherence condition it must not break, and a bounded attention budget.
We give Agent Smith a grammar, a type system, and a small-step lowering
semantics into agent objects, and we prove the compiler's central guarantee:
\emph{every well-typed Agent Smith compiles to an agent that honours, by
construction, four runtime invariants---conserved identity, a never-resetting
committed count, search-not-fetch world access, and construction/commitment
phase exclusion---and no ill-typed script is instantiated}. We prove the
language is \emph{mechanism-agnostic}: it specifies the agent's purpose and
the scenes that serve it, and delegates every domain behaviour to a
constructor-supplied hook, so that Agent Smith never contains the competence
of any domain (how to forge, how to bake, how to solve) but only the means to
instantiate an agent that pursues one. We prove a society of agents is
specified by the same language one level up, so any configuration of
agents---an ensemble instantiated and used as the specifier chooses---is
itself an Agent Smith. Because the agent an Agent Smith instantiates is, by
the underlying theory, indifferent between ``living a life'' and ``executing
a task''---the two are one object under two readings of where its purpose
rests---the language instantiates game characters and task-solving agents by
the same constructs. The compiler is the instantiator; the specifier
declares intent; the agent does the rest. Interpretive readings---of a script
as a character, of a society as a population---are confined to remarks on
which no theorem depends.
\end{abstract}

\vspace{0.5em}
\noindent\textbf{Keywords:} specification language; agent instantiation;
domain-specific language; typed compiler; lowering semantics; runtime
invariant; purpose-directed agent; non-player character; mechanism
agnosticism; society of agents.

% =====================================================================
\section{Introduction}
\label{sec:intro}
% =====================================================================

\subsection{Specifying an agent the way one specifies an NPC}

A designer who wants an iron smith in a game does not write the smith's
every hammer-blow. They say what the smith \emph{is}---a smith, with a
workshop, a standing in the town---and what the smith is \emph{for}: to
forge, to take contracts, to keep itself supplied with ore. The rest---which
of these the smith does at a given moment, how it forges a particular
blade---follows from what it is for and from the competence of forging, which
the designer does not re-derive. A convincing character is \emph{specified by
its purpose}, and its conduct is the pursuit of that purpose through whatever
activities serve it.

This paper gives a language for specifying agents in exactly this way, and a
compiler that instantiates them. The language is \textbf{Agent Smith}. A
user, an application, or another module writes an Agent Smith---a script that
declares an agent's purpose, the scenes in which it may act, and the few
standing quantities that make it an agent at all---and the compiler turns
that declaration into a running agent. The name is literal: every script
\emph{is} an Agent Smith, and the compiler is the smith that forges the agent
the script describes.

\subsection{What the language is, and is not}

Agent Smith is a \emph{specification and instantiation} language. Its scope
is deliberately narrow and its boundary is the crux of the design:

\begin{itemize}[leftmargin=*]
\item Agent Smith \emph{specifies} an agent: its purpose (what it is for),
its scenes (the activities that may serve that purpose), its identity
material, its coherence condition, its attention budget, its floors.
\item Agent Smith \emph{instantiates} an agent: the compiler reads a
well-typed script and produces a running agent that honours the runtime
invariants by construction.
\item Agent Smith does \emph{not} specify what the agent \emph{does} in a
domain. How a blade is forged, how bread is baked, how a task is solved is
\emph{not} in the language. That competence is supplied by the specifier
through a constructor hook (the declared drive and the per-scene behaviour),
exactly as the competence of an instrument is the player's and not the
score's. Agent Smith carries the intent; the mechanism is delegated.
\end{itemize}

The language thereby stays thin. It contains no domain grammar, no forging
routine, no baking routine, no solver; it contains only the constructs needed
to say \emph{what an agent is for} and to build one that pursues it.

\subsection{The underlying theory, used as a black box}

Agent Smith is a language \emph{over} a fixed theory of purpose-directed
agents, which it does not re-prove. That theory supplies, for any agent
satisfying a handful of axioms (finiteness, presence, costly individuation,
an act floor with bounded attention, phase exclusion): a conserved,
non-local, positive \emph{identity} invariant; a \emph{purpose} that is the
attractor of a strongly convex \emph{drive}; a \emph{water-filling} division
of attention across concurrent \emph{scenes} under a single price; a
\emph{monotone} committed history under which an agent never re-occupies a
past state and an authored copy is a distinct individual; \emph{search-not-%
fetch} world access, so a reply is walked afresh over the agent's present
graph rather than returned from a store; \emph{construction/commitment phase
exclusion}; \emph{two-factor relevance} (a response is admissible only if it
both advances purpose and preserves coherence); and, one level up, a
\emph{society} that is again an agent-shaped object with its own identity and
its own single shared purpose, coordinating in a common \emph{outcome space}
without any shared internals. Agent Smith's job is to let a specifier
\emph{name} the free parts of such an agent---its drive, its scenes, its
material---and to \emph{compile} the naming into an instance for which all of
the above holds. We treat the theory's results as given and cite them where
each is consumed; a reader who grants them can check every claim about the
language.

\subsection{Mechanism agnosticism, stated once}

The single most important property of Agent Smith is that it is
\emph{mechanism-agnostic}. The specifier says the agent is \emph{for} forging
(a drive whose attractor is ``blades exist, to specification'') and that
forging, contracting, and sourcing-ore are \emph{scenes} it may act in; the
specifier does \emph{not}, in Agent Smith, write how forging is performed.
That is supplied as a hook. Consequently the same language, with different
declared purposes and scenes, specifies a game smith and a task-solving agent
with no change to the language: by the underlying theory the two are one
object---an agent pursuing an attractor---differing only in whether the
attractor's cell is \emph{standing} (never saturated, so the agent runs on: a
character with a life) or \emph{attainable} (reached, so the agent halts: a
task done). Agent Smith instantiates both from the same constructs.

\subsection{Contributions}

\begin{enumerate}[leftmargin=*]
\item A grammar for Agent Smith (\Cref{sec:grammar}): the constructs by which
a user, application, or module specifies an agent by its purpose and scenes.
\item A denotation (\Cref{sec:denotation}): each construct is given a meaning
as an object of the underlying agent theory---self-graph, drive, scene,
coherence, budget, society.
\item A type system (\Cref{sec:typing}) that admits exactly the
specifications the compiler can instantiate into an agent honouring the
invariants, and a proof that typing is decidable.
\item A lowering semantics (\Cref{sec:lowering}) taking a well-typed script to
an agent instance, and the \emph{Instantiation Theorem}
(\Cref{sec:soundness}): every well-typed Agent Smith compiles to an agent that
honours the four runtime invariants by construction, and no ill-typed script
is instantiated.
\item The \emph{tick-loop} (\Cref{sec:tickloop}): a small-step specification of
what a compiled agent does---observe the solution-state (tiered, cheap-first),
diagnose whether it is stalled, and commit through a runtime-owned sufficiency
gate that reads the outcome, not the agent's reasoning---with a proof that the
loop preserves the four invariants and that a pure observer (which commits
nothing) is a valid, almost-free member.
\item The mechanism-agnosticism theorem (\Cref{sec:agnostic}): the language
carries no domain competence; behaviour is delegated to a constructor hook,
and the same script shape instantiates a character or a task-agent.
\item Society as an Agent Smith (\Cref{sec:society}): a configuration of
agents is specified by the same language one level up, closed under
composition.
\end{enumerate}

\subsection{Organisation}

\Cref{sec:objects} fixes the agent objects Agent Smith targets, imported from
the underlying theory. \Cref{sec:grammar} gives the grammar.
\Cref{sec:denotation} gives the denotation. \Cref{sec:typing} gives the type
system. \Cref{sec:lowering} gives the lowering semantics. \Cref{sec:tickloop}
specifies the tick-loop---what a compiled agent does each tick (observe,
diagnose, commit), with the pure observer as its floor case.
\Cref{sec:soundness} proves the Instantiation Theorem. \Cref{sec:agnostic}
proves mechanism agnosticism. \Cref{sec:lifecycle} covers management (spawn,
snapshot, retire). \Cref{sec:society} lifts to societies.
\Cref{sec:examples} gives worked Agent Smiths---a smith and a task-agent.
\Cref{sec:scope} states scope.

% =====================================================================
\section{The agent objects Agent Smith targets}
\label{sec:objects}
% =====================================================================

Agent Smith compiles specifications into \emph{agent objects}. We fix those
objects here, importing them from the underlying theory of purpose-directed
agents; each is used by the compiler as a black box with the stated
guarantee. Nothing in this section is proved anew.

\begin{definition}[Agent object]
\label{def:agent-object}
An \emph{agent object} is a tuple
\[
\Agent=\bigl(\Self,\ \drive,\ \scenes,\ \coh,\ \att,\ \floorc_a,\ m\bigr)
\]
where: $\Self=(\V,\E,\wt)$ is a finite connected weighted \emph{self-graph}
with weight floor $\floorc>0$ (the internal distinctions and the costs of
holding them apart); $\drive$ is a strongly convex \emph{drive potential} on
a compact convex disposition space $\mathcal{B}\subseteq\R^n$, with unique
minimiser the \emph{purpose} $\Purp=\arg\min\drive$; $\scenes$ is a finite
family of \emph{scenes}, each a pair $(A_\Scene,\gamma_\Scene)$ of an outcome
map $A_\Scene:\mathcal{B}\to\Out$ into a common outcome space $\Out$ and a
gain profile $\gamma_\Scene$; $\coh$ is a \emph{coherence} predicate on
candidate responses; $\att>0$ is the \emph{attention budget}; $\floorc_a>0$
is the \emph{act floor}; and $m\in\N$ is the \emph{committed count}.
\end{definition}

\begin{definition}[The four runtime invariants]
\label{def:invariants}
An agent object, once running, is required to preserve four
\emph{invariants}:
\begin{enumerate}[label=\textbf{(I\arabic*)},leftmargin=3em]
\item \emph{Conserved identity.} The character invariant
$\Char(\Agent)=\min_{\mathcal{Q}:k\ge2}\sum_{i<j}\wt(\cut(U_i,U_j))$ is
strictly positive, non-local, and unchanged under every relabelling of the
parts.
\item \emph{Never-resetting count.} $m$ is incremented on every committed act
and never decremented.
\item \emph{Search-not-fetch.} Every response is the terminus of a fresh
residual-descending walk on the present $\Self$; no answer is returned from a
store without a walk.
\item \emph{Phase exclusion.} At each instant the agent is in exactly one of a
construction phase (updating $\Self$, committing no act) or a commitment phase
(committing an act, forming no distinction).
\end{enumerate}
\end{definition}

\begin{principle}[The theory guarantees the invariants for any agent object]
\label{prin:theory}
For any agent object $\Agent$ (\Cref{def:agent-object}) whose self-graph has a
positive floor, whose drive is strongly convex, whose scenes have the stated
form, and whose acts are floored, the underlying theory guarantees:
\textup{(I1)} identity is a conserved positive non-local region;
\textup{(I2)} the committed count is strictly monotone and copies are
distinct; \textup{(I3)} recalling is searching; \textup{(I4)} construction and
commitment alternate; and, additionally, attention over the scenes divides by
a single-price water-filling rule, a response is admissible iff it is
two-factor relevant (purpose gain and coherence margin), and a society of
agent objects is again an agent object with its own identity and single shared
purpose, coordinating in $\Out$ without shared internals. We take these as
given.
\end{principle}

\begin{remark}[Agent Smith's whole job]
\label{rem:job}
\Cref{prin:theory} says the invariants hold for \emph{any} well-formed agent
object. Agent Smith's entire responsibility is therefore to guarantee that
what it compiles \emph{is} a well-formed agent object: a positive-floor
self-graph, a strongly convex drive, well-formed scenes, a floored act, a
coherence predicate, a positive budget. The type system (\Cref{sec:typing}) is
exactly the decidable check that a script denotes such an object; the
Instantiation Theorem (\Cref{sec:soundness}) is that a well-typed script
always does. The compiler does not re-prove the invariants; it earns them by
producing an object the theory already covers.
\end{remark}

% =====================================================================
\section{Grammar: how one specifies an agent}
\label{sec:grammar}
% =====================================================================

An Agent Smith is a script. We give its abstract grammar and then its concrete
syntax. The grammar has exactly the constructs needed to name the free parts
of an agent object (\Cref{def:agent-object}) and no others; in particular it
has no construct for a domain routine.

\subsection{Abstract grammar}

\begin{definition}[Agent Smith, abstractly]
\label{def:grammar}
An \emph{Agent Smith} $\spec$ is a term of the grammar
\[
\begin{array}{rcl}
\spec &::=& \texttt{agent}\ \iota\ \{\ \delta\ \sigma\ \rho\ \}
       \ \mid\ \texttt{society}\ \iota\ \{\ \overline{\spec}\ \theta\ \}\\[2pt]
\delta &::=& \texttt{purpose}\ \pi\quad(\text{the drive / attractor})\\[2pt]
\sigma &::=& \texttt{scenes}\ \{\ \overline{S}\ \}\\[2pt]
S &::=& \texttt{scene}\ \iota\ \texttt{serves}\ \pi\ \texttt{with}\ h\\[2pt]
\rho &::=& \texttt{self}\ \{\,\overline{p},\ \overline{e}\,\}\ \
           \texttt{budget}\ \att\ \ \texttt{floor}\ \floorc\ \
           \texttt{coherence}\ \coh\\[2pt]
\pi &::=& \texttt{minimise}\ \phi\ \mid\ \texttt{reach}\ o\in\Out\\[2pt]
\theta &::=& \overline{\texttt{tie}}\ \ \texttt{couple}\ K
\end{array}
\]
where $\iota$ is an identifier; $\phi$ names a strongly convex potential and
$o$ an outcome; $\overline{p},\overline{e}$ are declared parts and
separations; $h$ is a \emph{behaviour hook} (the delegated mechanism);
$\overline{\texttt{tie}}$ are inter-agent separations; $K$ is a coupling
strength. Overlines denote finite lists.
\end{definition}

\begin{remark}[There is no rule for ``what the scene does''---only that it
serves the purpose]
\label{rem:no-mechanism-rule}
The scene form is \texttt{scene}\ $\iota$\ \texttt{serves}\ $\pi$\
\texttt{with}\ $h$. It names the scene, declares \emph{which purpose it
serves}, and binds a hook $h$ that performs the domain work. The grammar has
no production expanding $h$ into forging steps or baking steps: $h$ is opaque
to Agent Smith, supplied by the specifier or a module. This is the
grammatical form of mechanism agnosticism (\Cref{sec:agnostic}): the language
can say a scene serves a purpose, never how it serves it.
\end{remark}

\subsection{Concrete syntax}

The concrete syntax is a lightweight surface over \Cref{def:grammar}. A
specification of a smith reads:

\begin{lstlisting}
// An Agent Smith: the iron smith.
// Says what the smith IS FOR; not how it forges.
agent smith {
  purpose minimise forge_residual
    // attractor: blades exist, to spec

  scenes {
    scene hammer   serves forge_residual with forge_hook
    scene contract serves forge_residual with contract_hook
    scene source   serves forge_residual with source_ore_hook
  }

  self {
    parts { skill, stock, orders, reputation }
    separations {
      (skill, stock: 3), (stock, orders: 2),
      (orders, reputation: 2), (reputation, skill: 2)
    }
  }
  budget 1.0
  floor  2.0
  coherence keeps { reputation, skill }
}
\end{lstlisting}

\noindent The three scenes---\texttt{hammer}, \texttt{contract},
\texttt{source}---are the smith's ``life'': all serve the one purpose
\texttt{forge\_residual}, and which the smith does at a moment is not scripted
but decided by water-filling over their gains (\Cref{prin:theory}). The hooks
\texttt{forge\_hook}, \texttt{contract\_hook}, \texttt{source\_ore\_hook}
carry the domain competence and are not part of Agent Smith.

\begin{remark}[The specifier writes intent, the compiler writes the agent]
\label{rem:intent}
Everything in the script above is \emph{intent}: what the smith is for, which
activities serve it, what holds it together, how much attention it has.
Nowhere does the script say ``at tick 5, hammer.'' The compiler
(\Cref{sec:lowering}) turns the intent into an agent object whose conduct is
the pursuit of \texttt{forge\_residual} through the declared scenes, and the
theory (\Cref{prin:theory}) supplies the conduct.
\end{remark}

% =====================================================================
\section{Denotation: what each construct means}
\label{sec:denotation}
% =====================================================================

We give each construct a denotation as a piece of an agent object
(\Cref{def:agent-object}). The denotation map $\llbracket\cdot\rrbracket$ is
the specification of the compiler's front end: it says what object a script is
asking for, before typing decides whether the request is admissible.

\begin{definition}[Denotation]
\label{def:denotation}
The denotation $\llbracket\spec\rrbracket$ of an Agent Smith is defined by
structural recursion:
\begin{align*}
\llbracket\texttt{self}\{\overline p,\overline e\}\rrbracket
  &= \Self=(\V,\E,\wt),\ \V=\overline p,\ \E=\overline e,\\
  &\qquad \wt(e)=\text{declared cost of }e;\\
\llbracket\texttt{purpose minimise }\phi\rrbracket
  &= \drive:=\phi,\ \Purp:=\arg\min\phi;\\
\llbracket\texttt{purpose reach }o\rrbracket
  &= \drive:=\tfrac12\,\dist(A(\cdot),o)^2,\ \Purp:=A^{-1}(o);\\
\llbracket\texttt{scene }\iota\texttt{ serves }\pi\texttt{ with }h\rrbracket
  &= \Scene=(A_\Scene,\gamma_\Scene),\ \mathrm{hook}(h);\\
\llbracket\texttt{budget }\att\rrbracket &= \att;\quad
\llbracket\texttt{floor }\floorc\rrbracket = \floorc;\quad
\llbracket\texttt{coherence }\coh\rrbracket = \coh;\\
\llbracket\texttt{agent }\iota\{\delta\sigma\rho\}\rrbracket
  &= \Agent=(\Self,\drive,\scenes,\coh,\att,\floorc_a,0);\\
\llbracket\texttt{society }\iota\{\overline\spec\,\theta\}\rrbracket
  &= \Soc=\text{quotient of }\llbracket\overline\spec\rrbracket\text{ under }\theta.
\end{align*}
The committed count of a freshly denoted agent is $0$ (I2's base).
\end{definition}

\begin{remark}[Two ways to name a purpose]
\label{rem:two-purposes}
\Cref{def:denotation} gives two forms of \texttt{purpose}. \texttt{minimise
}$\phi$ names the drive directly by a strongly convex potential---the general
case. \texttt{reach }$o$ names an outcome $o\in\Out$ and synthesises the drive
as squared distance to it---the common case, and the one that makes a
\emph{task}-agent: its attractor is the point $o$, attainable, so it halts at
quiescence. \texttt{minimise} with a floored potential whose minimum is a
standing region makes a \emph{character}: never saturated, so it runs on. The
same construct, two regimes (\Cref{thm:regimes}).
\end{remark}

% =====================================================================
\section{Type system: which specifications are instantiable}
\label{sec:typing}
% =====================================================================

Not every script denotes a \emph{well-formed} agent object. A drive that is
not strongly convex has no unique purpose; a self-graph with a zero-weight
separation has no floor; a scene serving no declared purpose is dead; a
non-positive budget admits no attention. The type system admits exactly the
scripts whose denotation is a well-formed agent object, so the compiler
instantiates all and only those.

\begin{definition}[Judgements]
\label{def:judgements}
The typing judgement $\judge\spec:\mathsf{Agent}$ reads ``$\spec$ is a
well-formed Agent Smith and denotes a well-formed agent object.'' Auxiliary
judgements type the parts: $\judge\rho:\mathsf{Self}$,
$\judge\delta:\mathsf{Drive}$, $\judge S:\mathsf{Scene}(\pi)$ (``scene $S$
serves purpose $\pi$'').
\end{definition}

\begin{trule}[Self well-formed]
\label{rule:self}
$\judge\texttt{self}\{\overline p,\overline e\}\ \texttt{floor}\ \floorc:
\mathsf{Self}$ holds iff the graph $(\overline p,\overline e)$ is connected,
simple, nonempty, and every declared cost is $\ge\floorc>0$. (Positive floor
and connectedness: the hypotheses of the identity theorem.)
\end{trule}

\begin{trule}[Drive well-formed]
\label{rule:drive}
$\judge\texttt{purpose }\pi:\mathsf{Drive}$ holds iff $\pi$ denotes a strongly
convex, continuously differentiable potential on a compact convex disposition
space (for \texttt{minimise }$\phi$, that $\phi$ is so; for \texttt{reach
}$o$, automatic, as squared distance is $1$-strongly convex). (Strong
convexity: the hypothesis of unique purpose and attraction.)
\end{trule}

\begin{trule}[Scene serves a declared purpose]
\label{rule:scene}
$\judge\texttt{scene }\iota\texttt{ serves }\pi\texttt{ with }h:
\mathsf{Scene}(\pi)$ holds iff $\pi$ is the agent's declared purpose, the gain
profile $\gamma_\Scene$ is nondecreasing with $\gamma_\Scene(0)=0$, and the
hook $h$ has the outcome-map type $\mathcal{B}\to\Out$. A scene serving no
declared purpose does not type. (No dead scenes: every scene must reduce the
one residual.)
\end{trule}

\begin{trule}[Agent well-formed]
\label{rule:agent}
\[
\frac{\judge\rho:\mathsf{Self}\quad \judge\delta:\mathsf{Drive}\quad
      \bigl(\judge S_i:\mathsf{Scene}(\pi)\bigr)_i\quad \att>0\quad
      \floorc_a>0}
     {\judge\texttt{agent }\iota\{\delta\,\texttt{scenes}\{\overline S\}\,
      \rho\}:\mathsf{Agent}}
\]
All parts well-formed, all scenes serving the one declared purpose, budget and
act floor positive.
\end{trule}

\begin{theorem}[Typing is decidable]
\label{thm:decidable}
Whether $\judge\spec:\mathsf{Agent}$ is decidable in time polynomial in the
size of $\spec$, given oracles for (a) connectivity of a finite graph and (b)
strong convexity of the named potential.
\end{theorem}

\begin{proof}
\Cref{rule:self} checks connectivity (linear in $\abs\E$), simplicity,
nonemptiness, and a floor comparison per edge---all polynomial.
\Cref{rule:drive} defers to the strong-convexity oracle; for \texttt{reach},
it is immediate. \Cref{rule:scene} checks a declared-purpose match (identifier
equality), monotonicity and the zero of the profile (given the profile's
form), and a hook type. \Cref{rule:agent} is a finite conjunction of the above
plus two positivity checks. Each premise is polynomial and there are finitely
many, so the judgement is decidable in polynomial time.
\end{proof}

\begin{remark}[The type system is exactly the well-formed-object check]
\label{rem:typing-is-wf}
\Cref{rule:self,rule:drive,rule:scene,rule:agent} are, one for one, the
hypotheses \Cref{prin:theory} needs: positive floor and connectedness
(identity), strong convexity (purpose), well-formed scenes (water-filling),
positive act floor (monotone count and search-not-fetch). Typing is not an
extra discipline layered on the theory; it \emph{is} the theory's hypotheses
made into a decidable front-end check.
\end{remark}

% =====================================================================
\section{Lowering: instantiation as compilation}
\label{sec:lowering}
% =====================================================================

The compiler $\compile$ takes a well-typed Agent Smith to a running agent
instance by a small-step lowering relation $\lowers$. Lowering is where
``specification'' becomes ``instance.''

\begin{definition}[Lowering]
\label{def:lowering}
The lowering relation $\spec\lowers\Agent$ is defined for
$\judge\spec:\mathsf{Agent}$ by:
\begin{enumerate}[label=\textup{L\arabic*.},leftmargin=2.6em]
\item \emph{Realise the self-graph.} Choose a concrete labelling of
$\llbracket\rho\rrbracket$; the identity invariant is label-independent (I1),
so any labelling is admissible.
\item \emph{Bind the drive.} Install $\drive=\llbracket\delta\rrbracket$ with
its minimiser $\Purp$; by \Cref{rule:drive} it is strongly convex, so $\Purp$
exists and attracts.
\item \emph{Install the scenes.} For each typed scene, register
$(A_\Scene,\gamma_\Scene)$ and bind its hook $h$; the compiler holds $h$ as an
opaque callable (\Cref{sec:agnostic}).
\item \emph{Set the standing quantities.} $\att,\floorc_a,\coh$ from $\rho$;
the committed count $m:=0$.
\item \emph{Seat the runtime.} Attach the water-filling scheduler over the
scenes, the two-factor relevance gate, the phase alternator, and the
search-walk responder---the fixed runtime of \Cref{prin:theory}.
\end{enumerate}
The result is an agent object $\Agent$ with $m=0$, ready to run.
\end{definition}

\begin{remark}[The compiler installs the runtime; it does not write conduct]
\label{rem:installs}
Step L5 attaches a \emph{fixed} runtime---the same scheduler, gate,
alternator, and responder for every agent, supplied once by the theory. The
compiler writes no conduct: the agent's behaviour is what that fixed runtime
produces when driven by \emph{this} agent's drive and scenes. Two different
Agent Smiths differ only in the objects installed at L1--L4; the machinery at
L5 is identical. This is why the language can be thin and the agents various.
\end{remark}

% =====================================================================
\section{The tick-loop: what a compiled agent does}
\label{sec:tickloop}
% =====================================================================

Lowering (\Cref{sec:lowering}) seats a fixed runtime at L5. We now specify
that runtime as a small-step relation---the \emph{tick-loop}---because it is
the one piece of an agent's behaviour that is common to every Agent Smith and
is worth pinning exactly. The loop realises the framework's governing stance:
\emph{an agent acts on the state of the solution, not on its own internal
reasoning}. Each tick, an agent reads the solution-state, diagnoses whether
it is stuck, and either commits an act or stays silent; nowhere does it
introspect a stored model of the whole task. The loop has three steps---%
\emph{observe}, \emph{diagnose}, \emph{commit}---and a distinguished floor
case, the pure observer, that commits nothing and costs almost nothing.

\subsection{The solution-state and its tiered reading}

\begin{definition}[Solution-state; residual gap]
\label{def:solution-state}
The \emph{solution-state} $\Omega$ is the shared, external record of the work
done so far toward the intent---the parts individuated, the boundaries
committed. It is \emph{not stored inside any agent}: no agent holds $\Omega$;
each reads a slice of it on demand. For an agent $\Agent$ with purpose
$\Purp$, the \emph{residual gap} $\Delta_\Agent(\Omega)\ge0$ is the semantic
distance from the current solution-state to $\Agent$'s attractor, floored (by
the underlying theory) at $\Delta_\Agent\ge\floorc>0$: the gap never closes to
a point.
\end{definition}

\begin{definition}[Reading tiers]
\label{def:tiers}
A \emph{reading} of $\Omega$ is tiered by cost:
\begin{itemize}[leftmargin=*]
\item \emph{Tier 0 (probe).} A deterministic slice retriever
$\mathrm{probe}(\Omega,q)$: given a query $q$ (what this agent needs), return
a ranked context slice of the present $\Omega$, cheaply and with no stored
answer---a re-derivation from $\Omega$ each call. (Concretely, an
empty-dictionary index over the substrate that returns a ranked slice per
query; it holds no domain answers, only the means to search for them.)
\item \emph{Tier 1 (deep read).} A model-backed extraction invoked only when a
Tier-0 slice does not suffice to move the residual---the escalation.
\end{itemize}
The agent always reads Tier 0 first and escalates to Tier 1 only on residue
(\Cref{def:step-observe}).
\end{definition}

\begin{remark}[Tier 0 is search-not-fetch, already]
\label{rem:tier0-search}
The probe of \Cref{def:tiers} is exactly the search-not-fetch discipline (I3)
made concrete: it returns a slice re-derived from the current $\Omega$, never
a stored answer, so a repeated probe against a grown $\Omega$ returns a
different slice. An agent that only ever reads Tier 0 therefore never fetches
and never holds the task; it re-reads the state. This is why the framework's
persistence is in the \emph{trajectory}---how the agent got here, its monotone
count and its position in the work---not in a cached task
(\Cref{rem:persistence}).
\end{remark}

\subsection{The three steps}

\begin{definition}[Observe]
\label{def:step-observe}
\emph{Observe} reads the solution-state for this agent's purpose:
\begin{enumerate}[label=\textup{O\arabic*.},leftmargin=2.6em]
\item probe Tier 0: $\varsigma\gets\mathrm{probe}(\Omega,q_\Agent)$;
\item measure the residual gap $\Delta_\Agent(\varsigma)$ from the slice;
\item if the slice suffices to determine the next act (the gap is legible from
Tier 0), stop; else escalate to Tier 1 and re-measure.
\end{enumerate}
Observe returns the current residual gap and a candidate next act, computed
from the solution-state slice alone---never from a stored model of the whole
task.
\end{definition}

\begin{definition}[Diagnose]
\label{def:step-diagnose}
\emph{Diagnose} reads the \emph{descent} of the residual gap across recent
ticks and returns a typed limit:
\begin{itemize}[leftmargin=*]
\item \textsc{converging}: the gap fell on the last committed act
($\Delta$ strictly decreased)---compute is the only obstacle; keep acting.
\item \textsc{compute-limited}: the gap is falling slowly but steadily---feed
more attention.
\item \textsc{structure-limited}: the gap has not fallen over a window (a
\emph{stall}) though acts were committed---more of the same act will not help;
the agent emits the stall as a typed signal so the society may route a
\emph{different} agent to the gap.
\end{itemize}
A stalled agent does not fall silent: it emits \textsc{structure-limited}, a
diagnosis the conductor layer reads.
\end{definition}

\begin{definition}[Commit; the sufficiency gate]
\label{def:step-commit}
\emph{Commit} is a \emph{runtime-owned} gate---not part of the agent's domain
work. It reads the \emph{outcome} of the candidate act (its effect on
$\Omega$ in the common outcome space) and fires iff two conditions hold:
\begin{enumerate}[label=\textup{(\roman*)},leftmargin=2.2em]
\item \emph{sufficiency}: the candidate act lowers the residual gap of the
\emph{parent} intent---its outcome-gain clears the attention price $\price$
(\emph{not} that the act reduces the agent's own gap to its floor: release at
sufficiency, not at completion);
\item \emph{coherence}: the act preserves the agent's coherence condition
(two-factor relevance).
\end{enumerate}
If both hold, the agent commits: it deposits an act, increments the count
$m$ (I2), and changes $\Omega$---which every other agent will read next tick.
If either fails, the agent commits nothing this tick and remains an observer.
\end{definition}

\begin{remark}[The gate reads the outcome, not the reasoning]
\label{rem:gate-outcome}
\Cref{def:step-commit} is the formal seat of ``act on the state of the
solution, not on internal reasoning.'' The commit decision is a function of
the act's \emph{outcome} in the shared space---does it move the parent's
residual, does it stay coherent---and never of how the agent's hook computed
the act. The domain work is opaque (\Cref{sec:agnostic}); the decision to act
is the runtime's, read off the outcome. This is why an agent needs no model of
any other agent: it reads the shared solution-state, not anyone's internals
(coordination in outcome space, \Cref{prin:theory}).
\end{remark}

\subsection{The loop, and its floor case}

\begin{definition}[Tick-loop]
\label{def:tick}
The \emph{tick} of a compiled agent $\Agent$ against solution-state $\Omega$
is
\[
\mathrm{tick}(\Agent,\Omega)\ =\ \mathrm{commit}\circ\mathrm{diagnose}\circ
\mathrm{observe}\,,
\]
in the construction/commitment phase order of I4: observe and diagnose run in
a construction phase (no act), commit runs in a commitment phase (at most one
act). The agent's behaviour is the iteration of $\mathrm{tick}$ against the
evolving $\Omega$.
\end{definition}

\begin{theorem}[The tick-loop preserves the four invariants]
\label{thm:tick-invariants}
Iterating $\mathrm{tick}$ (\Cref{def:tick}) preserves I1--I4. Identity is
untouched by any step (observe and diagnose read $\Omega$; commit changes
$\Omega$, not $\Agent$'s self-graph partition), so I1 holds; commit increments
$m$ and no step decrements it, so I2 holds; observe reads only via the probe,
which re-derives from the present $\Omega$ with no stored answer, so I3 holds;
and observe/diagnose occupy a construction phase while commit occupies a
commitment phase, disjoint by \Cref{def:tick}, so I4 holds.
\end{theorem}

\begin{proof}
By \Cref{def:step-observe,def:step-diagnose}, observe and diagnose compute
from a slice of $\Omega$ and the residual descent; they commit no act and do
not alter $\Self=(\V,\E,\wt)$, so the character invariant $\Char(\Agent)$
(\Cref{def:invariants}) is unchanged (I1) and the count is untouched. By
\Cref{def:step-commit}, commit deposits at most one act, incrementing $m$ by
one; no step lowers $m$ (I2). By \Cref{rem:tier0-search} the probe returns a
re-derived slice, never a stored answer, and Tier 1 is a model extraction
against the present state, so no response is fetched (I3). By \Cref{def:tick}
observe/diagnose are in construction and commit is in commitment, and I4's
phase exclusion makes these disjoint instants, so no tick both constructs and
commits (I4). Hence all four are preserved by every tick, and by induction by
the whole run.
\end{proof}

\begin{definition}[Pure observer; the floor case]
\label{def:observer}
An agent is a \emph{pure observer} at a tick if its commit gate does not fire:
either the candidate act's outcome-gain does not clear the price (the
solution-state does not yet call for it) or it would break coherence. A pure
observer runs observe (Tier 0 only---it never escalates, having no act to
justify the cost) and diagnose, commits nothing, and leaves $\Omega$ and its
own count $m$ unchanged.
\end{definition}

\begin{theorem}[The pure observer is valid and almost free]
\label{thm:observer}
A pure observer is a well-formed member of the society: it honours I1--I4
vacuously (it commits no act, so I2's count is unchanged and I4's commitment
phase is empty), it holds no task (I3: it only ever probes Tier 0), and it
costs only its Tier-0 probe per tick. It contributes nothing to the residual
($\delta=0$) yet remains present: the instant the solution-state changes so
that its candidate act clears the price, its commit gate fires and it acts. A
society may contain arbitrarily many pure observers at negligible cost.
\end{theorem}

\begin{proof}
A pure observer runs observe and diagnose (construction phase, no act) and a
commit gate that does not fire (\Cref{def:observer}); so it deposits no act,
leaving $m$ fixed (I2 vacuous) and the commitment phase empty (I4 vacuous),
while I1 (identity) and I3 (Tier-0 probe, re-derived slices) hold as for any
agent (\Cref{thm:tick-invariants}). Its per-tick cost is one Tier-0 probe, by
\Cref{def:observer}. Because the commit gate is re-evaluated each tick against
the current $\Omega$ (\Cref{def:step-commit}), a change in the solution-state
that raises the observer's candidate outcome-gain above the price fires the
gate on that tick: the observer acts exactly when the state calls for it, not
before. Hence it is present, valid, and almost free.
\end{proof}

\begin{remark}[The three at the zoo, and the crowd of cheap readers]
\label{rem:zoo}
\Cref{thm:observer} is the penguin-watcher: an agent attending to its own
scene, holding no model of the others, that commits the instant the shared
outcome-state changes (the others flee) because at that tick its
``get to safety'' act clears the price---coordination through the outcome, not
through anyone's internals. The same construction explains why a society
prefers \emph{many cheap readers} to one expensive one where a judgment is
uncertain: independent observers reading the same outcome-state sharpen the
collective judgment multiplicatively (\Cref{prin:theory}, crowd sharpening),
and each costs only its Tier-0 probe. The loop makes the cheap-crowd the
default and the deep read (Tier 1) the exception, paid only on escalation. We
use only the loop's step relation; the readings are orientation.
\end{remark}

% =====================================================================
\section{The Instantiation Theorem}
\label{sec:soundness}
% =====================================================================

We prove the compiler's central guarantee: well-typed scripts instantiate
invariant-honouring agents, and ill-typed scripts do not instantiate at all.

\begin{theorem}[Instantiation soundness]
\label{thm:instantiation}
If $\judge\spec:\mathsf{Agent}$ then $\spec\lowers\Agent$ for a unique agent
object $\Agent$ up to relabelling, and $\Agent$ honours the four runtime
invariants \textup{(I1--I4)} of \Cref{def:invariants} by construction.
\end{theorem}

\begin{proof}
By \Cref{rule:agent}, a well-typed $\spec$ has a well-formed self
(\Cref{rule:self}: connected, positive floor), a well-formed drive
(\Cref{rule:drive}: strongly convex), scenes each serving the declared purpose
(\Cref{rule:scene}), and positive $\att,\floorc_a$. Its denotation
$\llbracket\spec\rrbracket$ (\Cref{def:denotation}) is therefore a well-formed
agent object (\Cref{def:agent-object}). Lowering (\Cref{def:lowering})
realises it: L1--L4 produce exactly that object, and L5 attaches the fixed
runtime. By \Cref{prin:theory}, every well-formed agent object honours I1
(identity: positive floor $+$ connectedness give a conserved non-local
positive invariant), I2 (act floor gives a strictly monotone count), I3
(search-not-fetch: a floored act makes a zero-act fetch return nothing, so
every answer is a walk), and I4 (phase exclusion is part of the fixed
runtime). Uniqueness up to relabelling: L1's only freedom is the labelling,
under which the object is determined and, by I1, identity-invariant. Hence
$\Agent$ is unique up to relabelling and honours I1--I4.
\end{proof}

\begin{theorem}[No ill-typed instantiation]
\label{thm:no-illtyped}
If $\not\judge\spec:\mathsf{Agent}$ then $\compile$ emits no agent: the
compiler rejects $\spec$ at typing (\Cref{thm:decidable}) before lowering. In
particular no agent is ever instantiated with a zero-floor self, a non-convex
drive, a dead scene, or a non-positive budget.
\end{theorem}

\begin{proof}
$\compile$ runs the decidable type check (\Cref{thm:decidable}) before
lowering (\Cref{def:lowering}) and lowers only on
$\judge\spec:\mathsf{Agent}$. An ill-typed script fails some premise of
\Cref{rule:self,rule:drive,rule:scene,rule:agent}---the four failure modes
named---so it is rejected and never reaches $\lowers$. Hence no agent object
violating the theory's hypotheses is ever produced.
\end{proof}

\begin{corollary}[The compiler is the sole instantiator]
\label{cor:sole}
Every agent in the system is the image under $\compile$ of a well-typed Agent
Smith. There is no other route to an agent object, so every agent in the
system honours I1--I4 (\Cref{thm:instantiation}) and none violates the
hypotheses (\Cref{thm:no-illtyped}). Instantiation is centralised in the
compiler.
\end{corollary}

\begin{proof}
Agent objects arise only by lowering (\Cref{def:lowering}), which $\compile$
gates on typing. Combine \Cref{thm:instantiation,thm:no-illtyped}.
\end{proof}

% =====================================================================
\section{Mechanism agnosticism}
\label{sec:agnostic}
% =====================================================================

We prove the language's defining property: it carries no domain competence.
The behaviour of a scene is a constructor-supplied hook the compiler treats as
opaque, so the same language instantiates a game character and a task-agent
without change.

\begin{definition}[Behaviour hook]
\label{def:hook}
A \emph{behaviour hook} $h$ bound to a scene is a callable of type
$\mathcal{B}\times\text{Interaction}\to\mathcal{B}\times\Out$: given the
agent's disposition and an interaction, it returns the next disposition and an
outcome. Agent Smith holds $h$ as an opaque value: no production of
\Cref{def:grammar} inspects or expands $h$, and no typing rule constrains $h$
beyond its type (\Cref{rule:scene}).
\end{definition}

\begin{theorem}[The language contains no domain competence]
\label{thm:agnostic}
The denotation and typing of an Agent Smith are independent of the internal
behaviour of any hook: for two scripts $\spec,\spec'$ identical except that
each bound hook $h_i$ is replaced by a hook $h_i'$ of the same type, we have
$\judge\spec:\mathsf{Agent}\iff\judge\spec':\mathsf{Agent}$, and
$\llbracket\spec\rrbracket,\llbracket\spec'\rrbracket$ differ only in the
installed callables, not in $\Self,\drive,\att,\floorc_a,\coh$. Hence Agent
Smith specifies \emph{that} a scene serves the purpose, never \emph{how}.
\end{theorem}

\begin{proof}
Typing (\Cref{rule:self,rule:drive,rule:scene,rule:agent}) constrains a hook
only through its type $\mathcal{B}\to\Out$ (\Cref{rule:scene}); no rule reads
$h$'s body. Replacing $h_i$ by same-typed $h_i'$ leaves every premise
unchanged, so the judgements coincide. Denotation (\Cref{def:denotation})
records the hook as $\mathrm{hook}(h)$, an opaque binding; the self-graph,
drive, budget, floor, and coherence are computed from the other constructs and
are untouched by the swap. Hence the objects differ only in installed
callables.
\end{proof}

\begin{theorem}[One language, character or task-agent]
\label{thm:regimes}
Let $\spec_{\mathrm{char}}$ and $\spec_{\mathrm{task}}$ be Agent Smiths
identical except in their \texttt{purpose}: $\spec_{\mathrm{char}}$ uses
\texttt{minimise }$\phi$ with $\phi$'s minimum a standing region (never
saturated), and $\spec_{\mathrm{task}}$ uses \texttt{reach }$o$ with $o$ an
attainable outcome. Both type by the same rules and lower by the same
relation; the instantiated agents differ only in whether the attractor's cell
is standing or attained. Thus the same language, unchanged, instantiates a
character that lives on and a task-agent that halts.
\end{theorem}

\begin{proof}
Both \texttt{minimise }$\phi$ (with $\phi$ strongly convex) and \texttt{reach
}$o$ satisfy \Cref{rule:drive} (\Cref{rem:two-purposes}), so both scripts type
by \Cref{rule:agent} and lower by \Cref{def:lowering}. By \Cref{prin:theory}
the runtime is identical; the only difference is the drive installed at L2,
hence the location of $\Purp$. A standing-region minimum is never saturated
(the floor keeps residual positive), so the agent runs on; an attainable $o$
is reached at quiescence, so the agent halts. No construct other than
\texttt{purpose} differs. Hence one language instantiates both regimes.
\end{proof}

\begin{remark}[Why this is the whole point]
\label{rem:whole-point}
\Cref{thm:agnostic,thm:regimes} are the reason Agent Smith is a single small
language rather than a family of domain languages. Because it holds mechanism
as an opaque hook and purpose as the only thing distinguishing a character
from a task, a user specifies ``an iron smith'' and ``an agent that solves
this task'' with the same constructs, and the compiler instantiates each as an
agent pursuing its declared purpose through its declared scenes. The language
never learns to forge or to solve; it learns only to build an agent that is
\emph{for} forging or solving. The interpretation firewall of the underlying
theory is thereby inherited by the language: Agent Smith speaks only of
purpose, scenes, identity, budget---never of a domain.
\end{remark}

% =====================================================================
\section{Management: spawn, snapshot, retire}
\label{sec:lifecycle}
% =====================================================================

Agent Smith specifies \emph{and manages}. Beyond instantiation, the language
provides the lifecycle operations, each with a meaning that preserves the
invariants.

\begin{definition}[Lifecycle operations]
\label{def:lifecycle}
For a compiled agent $\Agent$:
\begin{itemize}[leftmargin=*]
\item \texttt{spawn}\ $\spec$: run $\compile(\spec)$ and register the resulting
$\Agent$ at count $0$.
\item \texttt{snapshot}\ $\Agent$: emit a serialisable record
$(\Self,\att,\floorc_a,\coh,\text{disposition},m,\text{phase})$---everything
but the opaque hooks, which are re-bound on rehydration.
\item \texttt{configure}\ $\Agent$: adjust budget $\att$ or coupling within the
running agent \emph{in the construction phase only} (I4).
\item \texttt{retire}\ $\Agent$: stop the agent; its snapshot persists, so a
later \texttt{spawn} from that snapshot resumes the \emph{same} individual at
its accumulated count (not a reset).
\end{itemize}
\end{definition}

\begin{theorem}[Lifecycle preserves the invariants]
\label{thm:lifecycle}
Each operation of \Cref{def:lifecycle} preserves I1--I4. In particular
\texttt{snapshot}/\texttt{spawn} round-trips the committed count, so a
rehydrated agent resumes at its accumulated $m$ and is the same individual;
and \texttt{configure} runs only in construction, so it never violates phase
exclusion.
\end{theorem}

\begin{proof}
\texttt{spawn} is $\compile$, covered by \Cref{thm:instantiation} (I1--I4 hold
at birth). \texttt{snapshot} records $m$ and the label-independent self-graph;
\texttt{spawn}-from-snapshot restores them, so $m$ is preserved (I2) and
identity is the same invariant (I1). Because the count is carried, the
rehydrated agent is not a fresh copy at $0$ but the same individual at $m$
(monotone-count distinctness of the underlying theory). \texttt{configure}
changes only $\att$ or coupling and only in the construction phase, so it
commits no act (I4 preserved) and does not touch the self-graph's partition
structure (I1 preserved) nor the count (I2). \texttt{retire} stops dispatch
and persists the snapshot; no count is lowered (I2). Search-not-fetch (I3) is a
property of the responder, attached at L5 and untouched by lifecycle. Hence all
four are preserved.
\end{proof}

\begin{remark}[Persistence is what makes the smith the same smith tomorrow]
\label{rem:persistence}
\Cref{thm:lifecycle} is the formal content of ``the smith you met yesterday,
mid-contract, is still that smith today.'' Because \texttt{snapshot}/%
\texttt{spawn} carries the committed count, resuming an agent is resuming the
\emph{same} individual with its accumulated history, not instantiating a
look-alike at count zero. A re-run of an agent's past act is therefore a new
act at a higher count---a genuine re-derivation against the grown graph, not a
replayed cache---which is exactly what a persistent character, or a resumable
task, requires.
\end{remark}

% =====================================================================
\section{A society is an Agent Smith}
\label{sec:society}
% =====================================================================

The language lifts to configurations of agents by the same constructs, one
level up. A \texttt{society} is itself an Agent Smith, so any configuration the
specifier wants is written in the same language.

\begin{definition}[Society specification]
\label{def:society-spec}
A \texttt{society}\ $\iota\ \{\ \overline\spec\ \theta\ \}$ denotes the
quotient object $\Soc$ whose vertices are the agents
$\llbracket\overline\spec\rrbracket$ and whose edges are the declared ties
$\theta$, with coupling $K$ for phase synchronisation. Its purpose, if
declared, is a society drive on the product disposition space whose minimiser
maps to the shared purpose in the common outcome space $\Out$.
\end{definition}

\begin{theorem}[Society well-formedness and closure]
\label{thm:society}
If each $\spec_i\in\overline\spec$ is well-typed and every declared tie has
cost $\ge\floorc$ and the society is connected, then
$\judge\texttt{society}\ \iota\{\overline\spec\,\theta\}:\mathsf{Agent}$: a
well-formed society is itself a well-formed agent object, honouring I1--I4 at
the society level, and coordinating in $\Out$ without shared internals. Hence
the language is closed under composition: an Agent Smith may be an agent or a
society of Agent Smiths.
\end{theorem}

\begin{proof}
By \Cref{prin:theory} a connected quotient of agent objects with
floor-$\ge\floorc$ ties is again an agent object (finite, connected, positive
floor), so \Cref{rule:agent} applies to $\Soc$: it has a conserved non-local
identity (I1 at agent level), a monotone society count (I2), search-not-fetch
and phase exclusion inherited by each member and by the quotient's own
scheduling (I3, I4). Coordination without shared internals is the
outcome-space coordination of the theory: members attain a shared region in
$\Out$ regardless of disjoint self-graphs. The society form is a production of
\Cref{def:grammar}, so the language is closed under it.
\end{proof}

\begin{remark}[Any configuration, written once]
\label{rem:configuration}
\Cref{thm:society} is what lets a specifier instantiate agents ``in whatever
configuration'': a lone smith, a smith and a miner tied by ore-supply, a whole
town, an ensemble of task-agents cooperating on a problem---each is a
\texttt{society} script, compiled by the same $\compile$, honouring the same
invariants. Because coordination lives in outcome space, the specifier need
give no agent a model of any other; declaring the ties and the shared outcome
suffices. The society's single shared purpose (one drive on the quotient) is
the formal sense in which the configuration acts as one.
\end{remark}

% =====================================================================
\section{Worked Agent Smiths}
\label{sec:examples}
% =====================================================================

We instantiate two agents from the same language: a character and a task-agent,
differing only in \texttt{purpose}.

\subsection{A character: the iron smith}

The smith of \Cref{sec:grammar}, whose purpose \texttt{minimise
forge\_residual} has a standing minimum: there is always more to forge, so the
agent runs on. Its three scenes are its life; water-filling picks which it
works at each moment; the necessity floor makes it ignore off-purpose requests
(a scene that does not serve \texttt{forge\_residual} never types, and an
off-purpose interaction gains below the attention price and is declined). It
persists across sessions by \texttt{snapshot}/\texttt{spawn}.

\subsection{A task-agent: rerun an experiment}

\begin{lstlisting}
// An Agent Smith: a task-solving agent.
// Same constructs; purpose is an attainable outcome.
agent rerun_exp {
  purpose reach verdict_confirmed
    // attractor: the experiment's verdict = CONFIRMED

  scenes {
    scene integrate serves verdict_confirmed with kuramoto_hook
    scene tabulate  serves verdict_confirmed with aggregate_hook
    scene report    serves verdict_confirmed with emit_hook
  }
  self {
    parts { data, method, result, verdict }
    separations {
      (data, method: 2), (method, result: 2),
      (result, verdict: 2), (verdict, data: 2)
    }
  }
  budget 1.0
  floor  2.0
  coherence keeps { method, result }
}
\end{lstlisting}

\noindent This types and lowers by the \emph{same} rules as the smith
(\Cref{thm:regimes}). Its purpose \texttt{reach verdict\_confirmed} is
attainable, so it halts at quiescence when the verdict is confirmed. A
\texttt{spawn} from a prior snapshot resumes the same task-individual at its
accumulated count, so ``redo the experiment'' re-derives against the grown
graph rather than replaying a stored result. The specifier wrote intent; the
compiler instantiated an agent that pursues it.

\begin{remark}[The same smith forges both]
\label{rem:same-smith}
The smith and the task-agent are two Agent Smiths, compiled by one compiler,
differing only in whether \texttt{purpose} names a standing region or an
attainable outcome. Everything else---the four invariants, the water-filling
over scenes, the persistence, the coordination if placed in a society---is
identical. This is the language's economy: to specify a game character and a
task-solver is to write the same kind of script, and to instantiate either is
one compilation.
\end{remark}

% =====================================================================
\section{Scope}
\label{sec:scope}
% =====================================================================

Agent Smith is a specification-and-instantiation language and nothing more. It
targets the agent objects of a fixed underlying theory, whose runtime
guarantees (the four invariants, water-filling attention, two-factor
relevance, outcome-space coordination) it consumes as a black box
(\Cref{prin:theory}) and does not re-prove. Its constructs name the free parts
of an agent---purpose, scenes, self, budget, floor, coherence, ties,
coupling---and its type system admits exactly the scripts whose denotation is a
well-formed agent object (\Cref{sec:typing}); its compiler instantiates all and
only those (\Cref{sec:soundness}). The language holds no domain competence:
scene behaviour is an opaque constructor hook (\Cref{sec:agnostic}), so the
same language instantiates a game character and a task-agent without change
(\Cref{thm:regimes}), and a configuration of agents is itself an Agent Smith
(\Cref{sec:society}). What the language does \emph{not} do is specify how any
scene performs its domain work, prescribe an agent's moment-to-moment conduct
(that is the runtime's, driven by the declared purpose and scenes), or
guarantee anything about a hook's internal correctness (that is the hook
author's). One hypothesis is inherited from the underlying theory and not
re-argued here: that an agent is a finite, positively-separated, floored,
purpose-directed object; absent it the invariants are not guaranteed and the
compiler's soundness does not hold. Interpretive readings---of a script as a
character with a life, of a society as a population, of a task-agent as an
instantiated task---are confined to remarks on which no theorem depends.

\subsection{Conclusion}

To specify a non-player character, one says what it is and what it is for and
lets its conduct follow; Agent Smith is the language for saying exactly that,
and its compiler is the instantiator that turns the saying into an agent.
Every script is an Agent Smith: a complete specification of one agent's
purpose, the scenes that serve it, and the material that makes it an agent at
all. The compiler types the specification against the hypotheses the
underlying theory requires, and lowers it into a running agent that honours,
by construction, a conserved identity, a never-resetting count,
search-not-fetch access, and phase exclusion---so that a smith is the same
smith tomorrow, an off-purpose request is ignored because attention costs
something, a copy is a new individual, and a re-run genuinely re-derives. The
language holds no domain competence: it says a scene serves the purpose, never
how, so the same constructs specify a smith who forges and an agent that
solves, a lone character and a whole society. A user, an application, or a
module specifies the agent the way a designer specifies an NPC; the compiler
forges it. That is the whole of Agent Smith: intent in, agent out, invariants
kept.

% =====================================================================
\bibliographystyle{plain}
\bibliography{references}
% =====================================================================

\end{document}
